\subsection{The politico-economic equilibrium with log}\label{\refsec:PEELOG}
Even with log preferences the political problem carries an endogenous state. Continuation policy is a function $\tau^{t+1}(\iota_t)$ of the informal savings ratio, and $\iota_t$ is shaped by the current tax through $\Theta_{s,t}$ in \eqref{\refeq:auxiliary:s0_s}. What we identify at each $t$ is therefore a policy \emph{function} of the predetermined $\iota_{t-1}$, not a number, and the tax path cannot be solved as one simultaneous system in $\bm{\tau}$: it has to be built up by a backward recursion over policy functions. As \eqref{\refeq:zdecomposition} and \eqref{\refeq:stateResidualLOG} below show, however, this particular state is a cheap one to carry, and the log case remains substantially less costly than the CRRA case of \cref{\refsec:PEE}.

Notationally, we let $z_t$ denote the marginal effect of bounded policies, for each $t$ defined as
\begin{align}\label{\refeq:focbounded}
    z_t &= \omega\sum_{j}\left(\gamma_{t-1,j}p_{t-1,j}\mu_{t-1,j} \frac{d\upsilon_{2,t}^j}{d\tau_t}\right)+\nu_t\sum\left(\gamma_{t,j}\mu_{t,j} \frac{d\upsilon_{1,t}^j}{d\tau_t}\right),
\end{align}
where the marginal effects on indirect utilities are defined by equations \eqref{\refmodeleq:PEELOG} for $t<T$ and \eqref{\refmodeleq:terminalPEELOG} for $t=T$.

\smalltitle{Grids.}
Let $\mathcal{T} = \left\lbrace \tau^{(1)},...,\tau^{(M)}\right\rbrace$ with $\tau^{(1)} = l$ and $\tau^{(M)} = u$ denote a grid of policy candidates on the interior, with $l\geq 0$ and $u<1$ strictly. The strict upper bound matters: both $d\upsilon_{1,t}^j/d\tau_t$ and $\partial\ln(\Theta_{h,t})/\partial\tau_t$ carry a factor $1/(1-\tau_t)$, so $z_t\rightarrow -\infty$ as $\tau_t\rightarrow 1$. Corners and multiplicity are handled by the selection rule \eqref{\refeq:candidates} of \cref{\refsec:RobustRoot}, applied along $\tau$ throughout this subsection.

Let $\mathcal{S}_0 = \left\lbrace\iota^{(1)},...,\iota^{(M_{\iota})}\right\rbrace$ with $\iota^{(1)} = l_{\iota}$ and $\iota^{(M_{\iota})} = u_{\iota}$ denote a grid of predetermined states, and let $\mathcal{S}_0'$ denote a separate grid of \emph{candidate} values of $\iota_t$. The two play different roles: $\mathcal{S}_0$ indexes the states we solve \emph{at}, while $\mathcal{S}_0'$ is searched \emph{over} to locate the state consistent with equilibrium.

\smalltitle{Align the candidate grid, do not refine it.}
The accuracy of the state search is governed by $\mathcal{S}_0'$, but it does not improve with refinement in the way one would expect, and the reason applies to any root located against an interpolated continuation policy. The functions $\tau^{t+1}(\cdot)$ and $\Theta_h^{t+1}(\cdot)$ are piecewise linear with breakpoints at the nodes of $\mathcal{S}_0$, and $\iota_t$ enters the residual \eqref{\refeq:stateResidualLOG} through nothing else. The residual inherits kinks at exactly those nodes and has no curvature of its own to speak of, so a cell of $\mathcal{S}_0'$ that \emph{straddles} a breakpoint is one across which the secant used to locate the crossing is not a faithful approximation. Taking $\mathcal{S}_0' = \mathcal{S}_0$ removes every straddling cell, and so does any refinement whose nodes \emph{contain} those of $\mathcal{S}_0$ (with error then falling as $O(h^2)$); a grid with more nodes but different ones performs materially worse and does not converge cleanly. The practical rule is to align rather than to refine, and to refine only in whole multiples if at all. The companion residual in the savings level (\cref{\refsec:PEE}) behaves in the opposite way, for a reason given there that makes the contrast instructive rather than confusing.

\smalltitle{The state grid is strictly positive.}
We take $l_{\iota}>0$ strictly, which is not innocuous and is worth being explicit about. From \eqref{\refmodeleq:EE:c0}, old-age informal consumption is proportional to $\iota_{t-1}+A_t\tau_t$, and the corresponding term of the first order condition,
\begin{align}\label{\refeq:dv20}
    \frac{d\upsilon_{2,t}^0}{d\tau_t} = (1-\alpha)\frac{d\ln\left(h_t\right)}{d\tau_t}+\frac{A_t}{\iota_{t-1}+A_t\tau_t},
\end{align}
has a pole where that bracket vanishes. Since $A_t>0$ and $\tau_t\geq 0$, requiring $\iota_{t-1}>0$ places the pole outside the evaluation region for every node of $\mathcal{T}$, so $z_t$ is bounded on the whole grid and no masking of the state dimension is needed. What the restriction assumes is that the informal household is a net saver --- see the footnote to \eqref{\refmodeleq:informalBudget} --- and that assumption is verified against the solved path rather than taken for granted: a solved $\iota_t$ that falls outside $[l_{\iota},u_{\iota}]$ is treated as an infeasibility and reported as such; it is never clipped back onto the grid.

\smalltitle{Bounds for the state grid.}
Bounds are available from the steady state, though not by reading its range off directly. Under a constant policy $\tau$, \eqref{\refeq:auxiliary:s0_s} evaluated at the steady-state $\Theta_s(\tau)$ of \cref{\refsec:EE} gives a steady-state ratio $\iota^*(\tau)$, and two of its features rule out the naive construction. First, $\iota^*(\tau)\rightarrow\infty$ as $\tau\rightarrow 1$ --- not because informal savings explode, but because the denominator of $\iota$ is \emph{formal} savings, which the tax crowds out entirely --- so the upper end of $\iota^*$'s range describes a degenerate corner that no equilibrium path approaches, and a grid spanning it would place essentially all of its nodes where no path can go. Second, $\iota^*$ \emph{understates} the range $\iota_t$ actually takes along the finite-horizon recursion, particularly near the terminal period, so a lower bound padded from $\iota^*$'s own minimum alone leaves part of $\mathcal{T}$ infeasible for a reason with no economic content, and the political selection is then pinned to the edge of the feasible set rather than determined by the first order condition.

What \emph{is} a reliable anchor is $\min_{\tau}\iota^*(\tau)$: it is attained in the interior, finite by construction, and stable across the intertemporal elasticity and across preference specifications. The set of states the recursion actually reaches sits at stable multiples of it, and we take
\begin{align*}
    l_{\iota} = \max\left\lbrace\delta_{\iota},\, \underline{\lambda}\min_{\tau}\iota^*(\tau)\right\rbrace, \qquad
    u_{\iota} = \min\left\lbrace\overline{\lambda}\min_{\tau}\iota^*(\tau),\, \bar{u}\right\rbrace,
\end{align*}
with $\underline{\lambda} = 0.45$, $\overline{\lambda} = 3.7$, a small $\delta_{\iota}>0$, and nodes spaced logarithmically --- well defined precisely because $l_{\iota}>0$. The multiples carry roughly $20\%$ margin on the reachable set's measured edges. The absolute cap $\bar{u}$ is a backstop rather than the operative bound; if it ever binds, that signals the anchor has moved and the multiples should be re-measured, not the cap re-chosen.

Two diagnostics govern any adjustment of these bounds. When tightening the lower bound, the quantity to watch is the number of \emph{corner selections}, not the number of feasible nodes: nodes lost at the top of $\mathcal{T}$ (where the implied $\iota_t$ exceeds a bound the equilibrium never approaches) are harmless, while nodes lost at the bottom pin the selection, and raw feasibility counts do not distinguish the two. And the two failure directions are asymmetric: bounds that are too narrow announce themselves, since a solved $\iota_t$ outside the grid is reported as an infeasibility, whereas bounds that are too wide are silent --- they cost resolution where the dynamics live and nothing warns of it. The cheap detector for the second failure is the reachable set: reporting the fraction of $\mathcal{S}_0$'s nodes that fall inside it alongside each solve turns a silent inefficiency into a number one can read (around $0.8$ under the rule above).

\smalltitle{Why $s_{t-1}$ is not a state, and $\iota_{t-1}$ is.}
By \eqref{\refmodeleq:EE:sigma_ci}, every equilibrium consumption function takes the form $x_t = \Theta_{x,t}\left(s_{t-1}/\nu_t\right)^{\sigma_x}$ with an exponent $\sigma_x$ that does not depend on policy. Taking logs, $\ln(x_t)$ is affine in $\ln(s_{t-1})$ with a policy-independent slope, so
\begin{align}\label{\refeq:logsep}
    \frac{d\ln\left(x_t\right)}{d\tau_t} = \frac{d\ln\left(\Theta_{x,t}\right)}{d\tau_t},
\end{align}
and the same holds for $h_t$ by \eqref{\refeq:auxiliary:h} and, after substituting \eqref{\refeq:auxiliary:h} and \eqref{\refeq:auxiliary:s} into \eqref{\refeq:auxiliary:R}, for $R_{t+1}$. Under log preferences the marginal indirect utilities are pure log-derivatives, so \emph{no} term of $z_t$ depends on the savings level $s_{t-1}$: it is not a state of the political problem, and the entire first order condition can be evaluated from the coefficient functions $\Theta$ alone.

What survives are the two predetermined ratios that the old generations carry into $t$. Of these, $s_{t-1,i}/s_{t-1}$ is a closed-form function of $\tau_t$ through \eqref{\refeq:auxiliary:si_s} evaluated at the previous period's vintage, so it is not a state either; this is the reduction recorded in \cref{\refmodelsec:PEE}, and it is what keeps the whole distribution of past savings out of the problem. The informal ratio $\iota_{t-1}$ admits no such reduction: by \eqref{\refeq:auxiliary:s0_s} it depends on $\Theta_{s,t-1}$, hence on $\tau_{t-1}$, which the policy maker at $t$ takes as given. It is the single state of the problem.

\smalltitle{The state enters through one additive term.}
By \eqref{\refeq:dv20}, $\iota_{t-1}$ reaches $z_t$ only through the old informal generation's term, and only through the scalar combination $\iota_{t-1}+A_t\tau_t$. Hence
\begin{align}\label{\refeq:zdecomposition}
    z_t\left(\tau_t,\iota_{t-1}\right) = \bar{z}_t\left(\tau_t\right)+\omega\gamma_{t-1,0}p_{t-1,0}\mu_{t-1,0}\frac{A_t}{\iota_{t-1}+A_t\tau_t},
\end{align}
where $\bar{z}_t$ collects every other term of \eqref{\refeq:focbounded} --- including the $(1-\alpha)d\ln(h_t)/d\tau_t$ part of \eqref{\refeq:dv20} --- and is a function of $\tau_t$ alone. Evaluating $z_t$ on all of $\mathcal{T}\times\mathcal{S}_0$ therefore costs one pass over $\mathcal{T}$ plus a rank-one correction. Refining $\mathcal{S}_0$ is close to free, which is why a state grid is affordable here at a resolution that would be uncomfortable for a state entering the expensive part of the problem.

\smalltitle{The informal savings ratio is a fixed point in $\iota_t$.}
Suppose period $t+1$ is solved and we hold the continuation policy $\tau^{t+1}(\cdot)$ together with the coefficient function $\Theta_h^{t+1}(\cdot)$, both functions of the state $\iota_t$ entering $t+1$. Note that we carry $\Theta_{h,t+1}$ rather than $h_{t+1}$: by \eqref{\refeq:auxiliary:h} the level $h_{t+1}$ depends on $s_t$ as well as on $\iota_t$ and so is not a function of the state, whereas its logarithm differs from $\ln(\Theta_{h,t+1})$ by a term that is constant along $\tau_t$, which by \eqref{\refeq:logsep} is all the first order condition needs.

For a candidate current tax $\tau_t$ and a candidate $\iota_t$, the auxiliary equations \eqref{\refeq:auxiliary} evaluate in one forward pass, using $B_{t+1}^i = \beta_{t,i}$ and $B_{t+1}^0 = \beta_{t,0}$:
\begin{subequations}\label{\refeq:stateApproxLOG}
\begin{align}
    \Gamma_{s,t} &= \Gamma_{s,t}\left(\bm{\beta}_t,\, \tau^{t+1}(\iota_t)\right), \label{\refeq:stateApproxLOG:Gammas} \\
    \Theta_{h,t} &= \Theta_{h,t}\left(\tau_t,\, \tau^{t+1}(\iota_t),\, \theta_{t+1},\, \Gamma_{s,t}\right), \label{\refeq:stateApproxLOG:Thetah} \\
    \Theta_{s,t} &= \Theta_{s,t}\left(\Theta_{h,t},\, \Gamma_{s,t}\right). \label{\refeq:stateApproxLOG:Thetas}
\end{align}
\end{subequations}
The candidate is an equilibrium precisely when it reproduces itself through \eqref{\refeq:auxiliary:s0_s}. Defining the residual
\begin{align}\label{\refeq:stateResidualLOG}
    \mathcal{R}_t^0\left(\tau_t,\iota_t\right) \equiv \iota_t\left(\beta_{t,0},\, \Theta_{s,t}\left(\tau_t,\iota_t\right),\, \tau^{t+1}(\iota_t)\right)-\iota_t,
\end{align}
the state policy function $\iota_t(\tau_t)$ is the root of \eqref{\refeq:stateResidualLOG} in its second argument. This is a genuine root-finding problem --- $\iota_t$ enters the right-hand side through the continuation policy --- and it is a \emph{root} problem, not a maximisation: any crossing is a solution, so the direction-filtering and objective-integration apparatus of \eqref{\refeq:candidates} is not used here.

\smalltitle{The fixed point does not involve the state.}
Read \eqref{\refeq:stateApproxLOG} and \eqref{\refeq:stateResidualLOG} again and note that $\iota_{t-1}$ appears nowhere in them. This is a stronger statement than it may look, and it is worth contrasting with the CRRA case of \cref{\refsec:PEE}: there the predetermined savings level survives in the corresponding residual as an explicit factor $\left(s_{t-1}/\nu_t\right)^{\alpha(1+\xi)/(1+\alpha\xi)}$, so the residual has to be broadcast across the state grid before any root is located. Here the state drops out of the fixed point entirely: with $B_{t+1}^i = \beta_{t,i}$ a primitive, $\Gamma_{s,t}$ is a function of $\tau_{t+1}$ alone, and $\Theta_{s,t}$ is a coefficient function in which the level of past savings has already been divided out. The state approximation is therefore genuinely one-dimensional --- solved once over $\mathcal{T}\times\mathcal{S}_0'$, yielding $\iota_t(\tau_t)$ on $\mathcal{T}$ --- and is not repeated per state. Together with \eqref{\refeq:zdecomposition}, a period of the recursion costs one pass over $\mathcal{T}\times\mathcal{S}_0'$ and two broadcasts.

\smalltitle{Feasibility.}
The root of \eqref{\refeq:stateResidualLOG} need not lie inside $\mathcal{S}_0'$: at large $\tau_t$ the coefficient $\Theta_{s,t}$ is depressed enough that the implied $\iota_t$ leaves the grid from above, and the discounted-pension term can push it out from below. Outside that window no state is located and every object downstream of $\iota_t$ is undefined. We therefore carry an explicit feasibility mask, marking the $\tau_t\in\mathcal{T}$ for which a root was found, and require at least two feasible nodes before attempting to select a maximum. By the previous paragraph this mask is one-dimensional --- it is the same for every $\iota_{t-1}$ --- which is a direct consequence of the state not entering the fixed point, and is verified numerically rather than assumed.

\smalltitle{Which derivatives are taken on the grid.}
From \eqref{\refmodeleq:PEELOG}, the young generations' indirect utilities are
\begin{align}\label{\refeq:v1LOG}
    \upsilon_{1,t}^i = \left(1+\beta_{t,i}\right)\ln\left(\tilde{c}_{1,t}^i\right)+\beta_{t,i}\ln\left(R_{t+1}\right)+\text{const}, \qquad
    \upsilon_{1,t}^0 = \left(1+\beta_{t,0}\right)\ln\left(\tilde{c}_{1,t}^0\right)+\beta_{t,0}\ln\left(R_{t+1}\right)+\text{const},
\end{align}
where the constants involve no policy, and where the informal household's own return enters as $\ln(R_{t+1}^0) = \ln(R_{t+1})+\ln(\chi_{t+1}^R)$ with the second term exogenous, so both lines carry the same $R_{t+1}$. Both depend on $\tau_t$ through the continuation policy $\tau^{t+1}\left(\iota_t(\tau_t)\right)$, which is an interpolant with no closed-form derivative. These terms are therefore \emph{not} available in closed form --- the presence of the state is what costs us that --- and must be differentiated numerically along $\tau_t$. It is convenient to assemble each profile in \eqref{\refeq:v1LOG} first and differentiate it in one step, rather than differentiating $\ln(\tilde{c}_{1,t}^j)$ and $\ln(R_{t+1})$ separately. By \eqref{\refeq:logsep} the profiles need only be known up to an additive constant, so they may be built from the $\Theta$'s without reference to $s_{t-1}$ at all.

The old generations' terms are not differentiated on the grid:
\begin{itemize}
    \item $d\upsilon_{2,t}^i/d\tau_t$ \textbf{must} be taken from its closed form in \eqref{\refmodeleq:PEELOG}. Differentiating it numerically along the grid would be wrong, not merely imprecise: $s_{t-1,i}/s_{t-1}$ varies along that grid, and the policy maker takes it as predetermined, so a grid derivative would include a channel that does not belong in the first order condition.
    \item $d\upsilon_{2,t}^0/d\tau_t$ carries no such hazard, and this deserves emphasis because it is the opposite of what the analogy with $s_{t-1,i}/s_{t-1}$ would suggest. $\iota_{t-1}$ is a grid coordinate, held fixed by construction as $\tau_t$ varies, so a grid derivative would hold it fixed correctly. Promoting the ratio to a state is precisely what removes the hazard. We nonetheless use the closed form \eqref{\refeq:dv20}, which is exact and cheaper.
\end{itemize}
Both closed forms require $d\ln(h_t)/d\tau_t$, which is itself numerical here and is obtained as $d\ln(\Theta_{h,t})/d\tau_t$ along $\mathcal{T}$, by \eqref{\refeq:logsep}.

\smalltitle{The coordinate the numerical derivatives are taken in.}
One detail of that differentiation is not cosmetic. Every profile above carries a $\ln(1-\tau_t)$ term, through $\Theta_{h,t}\propto\left((1-\alpha)(1-\tau_t)\right)^{\xi/(1+\alpha\xi)}$, so each derivative diverges like $1/(1-\tau_t)$ as $\tau_t\rightarrow u$. An interpolant fitted through uniformly spaced nodes of $\mathcal{T}$ cannot represent that, and the error concentrates exactly where the first order condition is steepest. Fitting instead in the coordinate
\begin{align*}
    x = \ln\left(1-\tau_t\right), \qquad\text{so that}\qquad \frac{dy}{d\tau_t} = \frac{dy}{dx}\cdot\frac{-1}{1-\tau_t},
\end{align*}
makes the singular part affine in $x$. It is then reproduced exactly, and only the genuine curvature coming from $\tau^{t+1}\left(\iota_t(\tau_t)\right)$ is left to approximate; measured against the closed form available when the continuation policy is held fixed, the substitution improves the derivative from percent-level error near the top of $\mathcal{T}$ to machine precision. The same substitution is used in \cref{\refsec:PEE}. What remains after it is the kink structure inherited from the interpolated continuation policy, which sets a small floor on the accuracy of $z_t$; near a root this displaces the located tax by a small fraction of the spacing of $\mathcal{T}$, so the grid, not the differentiation, is what limits the solved policy.

\begin{alg}[Backward grid search over the informal savings ratio]
\label{\refalg:LOG:gridsearch}
Fix $\bm{\theta},\bm{\epsilon}$ and the grids $\mathcal{T}$, $\mathcal{S}_0$, $\mathcal{S}_0'$. The recursion identifies a sequence of policy functions $\tau_t(\iota_{t-1})$ together with the state policy functions $\iota_t(\tau_t)$ that determine next period's state, iterating backwards from $t=T$.

\smallskip\noindent\emph{Terminal period.} There is no continuation policy, hence no fixed point to solve: by \cref{\refmodelsec:finitehorizon} every object entering $z_T$ is closed form in $(\tau_T,\iota_{T-1})$. On $\mathcal{T}$ evaluate $\Theta_{h,T}(\tau_T)$ from \eqref{\refeq:auxiliary:ThetahT} and $d\ln(h_T)/d\tau_T = -\frac{\xi}{1+\alpha\xi}\frac{1}{1-\tau_T}$, then $\bm{B}_T$, $\Gamma_{s,T-1}$ and $s_{T-1,i}/s_{T-1}$ at the previous period's vintage. The young terms are closed form by \eqref{\refmodeleq:terminalPEELOG}, with $d\upsilon_{1,T}^0/d\tau_T = 0$ --- the informal young neither save nor face a future at $T$ --- so no numerical differentiation is needed in the terminal period at all. Assemble $\bar{z}_T$ on $\mathcal{T}$, broadcast to $\mathcal{T}\times\mathcal{S}_0$ through \eqref{\refeq:zdecomposition}, and apply \eqref{\refeq:candidates} along $\tau$ for each $\iota_{T-1}$ to obtain $\tau_T(\iota_{T-1})$. Finally evaluate $\Theta_{h,T}$ at the selected taxes and construct the interpolants $\tau^T(\cdot)$ and $\Theta_h^T(\cdot)$ on $\mathcal{S}_0$.

\smallskip\noindent\emph{Recursion, $t = T-1,...,1$.} Given $\tau^{t+1}(\cdot)$ and $\Theta_h^{t+1}(\cdot)$, each period proceeds in four steps.
\begin{enumerate}
    \item \emph{State approximation.} Over $\mathcal{T}\times\mathcal{S}_0'$ evaluate the forward pass \eqref{\refeq:stateApproxLOG} and locate the root of \eqref{\refeq:stateResidualLOG} along the $\mathcal{S}_0'$ dimension, giving $\iota_t(\tau_t)$ on $\mathcal{T}$ together with its feasibility mask. This is the only step that touches $\mathcal{S}_0'$, and the only one that calls the continuation interpolants.
    \item \emph{Current-period objects.} On $\mathcal{T}$, with $\tau_{t+1} = \tau^{t+1}\left(\iota_t(\tau_t)\right)$ and $\Theta_{h,t+1} = \Theta_h^{t+1}\left(\iota_t(\tau_t)\right)$, evaluate $\Gamma_{s,t}$, $\Theta_{h,t}$, $\Theta_{s,t}$, and the two young profiles of \eqref{\refeq:v1LOG} up to their $s_{t-1}$-constants. Evaluate $\bm{B}_t$, $\Gamma_{s,t-1}$ and $s_{t-1,i}/s_{t-1}$ at the previous period's vintage, exactly as in the terminal period.
    \item \emph{Derivatives and the first order condition.} Differentiate the two young profiles numerically along $\tau_t$, in the coordinate $x = \ln(1-\tau_t)$; obtain $d\ln(h_t)/d\tau_t$ as $d\ln(\Theta_{h,t})/d\tau_t$ the same way; take $d\upsilon_{2,t}^i/d\tau_t$ and the first term of \eqref{\refeq:dv20} from their closed forms. Assemble $\bar{z}_t$ on $\mathcal{T}$, leaving infeasible nodes undefined, and broadcast to $\mathcal{T}\times\mathcal{S}_0$ through \eqref{\refeq:zdecomposition}.
    \item \emph{Selection and reporting.} Apply \eqref{\refeq:candidates} along $\tau$ for each $\iota_{t-1}$, skipping infeasible nodes, to obtain $\tau_t(\iota_{t-1})$. Re-evaluate steps 1--2 at the selected taxes to obtain the solution on $\mathcal{S}_0$, and construct the interpolants $\tau^t(\cdot)$ and $\Theta_h^t(\cdot)$ that period $t-1$ will call. The selected $\iota_t$ may be recorded as an interpolant over $\iota_{t-1}$ as well; period $t-1$ never calls it, and it serves only the cheaper variant of the path solve in \cref{\refsec:PEEpath}. A light smoothing of $\tau_t(\iota_{t-1})$ before interpolation removes small kinks inherited from the interpolated continuation policy, subject to three requirements. The smoothed values are clipped back into $[l,u]$: a smoother run through a profile that is flat at a corner over part of $\mathcal{S}_0$ will undershoot it, and the reported tax would otherwise leave the admissible set wherever many states select a corner. The smoothing is confined to this reporting step and never applied to the interpolant through which step 3's derivatives are taken --- in the coordinate used there the profiles are already smooth, so there is nothing to denoise, and a smoothing budget large enough to matter degrades the derivative rather than improving it. And the smoother's interior knots are \emph{fixed} in advance --- at every $m$-th valid node of the grid being smoothed, placed on the surviving nodes where a profile carries infeasible cells, so that each knot interval retains data --- rather than selected from the data by a residual criterion. Fixed knots make the smoothed profile a \emph{linear} map of the profile it smooths, hence continuous in the model parameters; a smoother that chooses its own knot count introduces a discrete choice, and hence jumps, inside a function that the calibration of \cref{\refsec:calibration} is about to differentiate, at which point an outer root lying inside a jump simply does not exist in the discretised problem.
\end{enumerate}
The cost is $O\left(T\cdot M\cdot M_{\iota}'\right)$ evaluations of closed-form expressions, with no initial guess for the path. Note that the state grid $\mathcal{S}_0$ enters this count only through the selection step, not through the evaluations, which is the practical content of \eqref{\refeq:zdecomposition}.
\end{alg}
